\SourceLeaf{169}
\Needspace{9\baselineskip}
\pdfbookmark[1]{[IV.] Bilanz zwischen den Verpflichtungen und den Mitteln der Professoren-Witwenkasse zu Göttingen.}{witwen-iv}
\begin{center}
[IV.]\par\medskip
\textit{Bilanz zwischen den Verpflichtungen und den Mitteln der Professoren-Witwenkasse zu Göttingen.}
\end{center}
\noindent\hspace*{\parindent}In Folge des von der Universitäts-Kirchendeputation vor einigen Monaten mir eröffneten Wunsches habe ich mich der neuen Berechnung der Bilanz der Professoren-Witwenkasse unterzogen und diese Arbeit jetzt vollendet. Ich werde meinen Bericht darüber so anordnen, dass ich zuerst die nöthigen allgemeinen Erläuterungen vorausschicke; hiernächst die Bilanzrechnung selbst in einer concisen leicht übersichtlichen Form aufstelle; sodann die einzelnen Posten der Rechnung näher erortere, und endlich die Vorschläge fur die nächst bevorstehende Periode daran knupfe.
\begin{center}\rule{.26\textwidth}{.3pt}\end{center}
\noindent\hspace*{\parindent}Die Bilanzaufstellung für ein solches Institut wie unsre Witwenkasse muss sich offenbar auf einen bestimmten \textit{Zeitpunkt} beziehen. Dass ich dafür diesmal den 1. October 1851 gewählt habe, wird keiner weitläuftigen Rechtfertigung bedürfen. Die erste Bilanzrechnung war für den 1. October 1845 gestellt gewesen: aus nahe liegenden Gründen muss die Zwischenzeit zwischen zwei auf einander folgenden Prufungen eine volle Anzahl von Jahren umfassen, um das Resultat reiner hervortreten zu lassen: endlich, wenn in Folge der neuen Prüfung eine Modification der bisherigen Punktationen als angemessen erscheinen sollte, so wird man bei dem Beschluss offenbar viel lieber sich auf den neuesten Zustand stutzen wollen, als auf denjenigen, welcher vor einem Jahre Statt gefunden hat. Die jetzigen Statuten schreiben zwar aller funf Jahre eine neue Revision vor: allein der Zeitpunkt, wo die Aufforderung an mich gelangte, liess eine andere Wahl nicht mehr zu; auch ist durch diese Erstreckung des funfjährigen Zeitraumes auf einen sechsjährigen nicht nur nichts verloren, sondern vielmehr eine noch etwas entschiedenere Auspragung der Zustandsänderung gewonnen.

Das Wesen der ganzen Bilanzrechnung der jetzigen wie der von 1845, besteht darin, dass nicht fur das nächste Jahr und nicht fur einige Jahre, sondern fur alle Zukunft, einerseits die Obliegenheiten des Instituts, andererseits seine Hülfsmittel auf den äquivalirenden Capitalwerth zurückgeführt werden. Nur auf diesem Wege ist es möglich, einer Anstalt, die nur zum kleinsten Theile auf Beiträge, und dem grössten Theile nach auf ihren Vermogensbesitz basirt, und in den letzten Decennien an Theilnehmerzahl so sehr vergrössert ist, die Haltbarkeit für alle Zukunft zu sichern.

Obgleich diesmal eben so wie 1845 alle Rechnungen doppelt geführt sind, nemlich nach dem Zinsfuss von \(3\tfrac{1}{2}\) und nach dem von 4 Procent, so habe ich es doch fur hinreichend gehalten, hier nur die Resultate nach ersterem aufzuführen. Ein Theil des Vermögens trägt wirklich nur \(3\tfrac{1}{2}\) Procent; von einem andern jetzt hoher verzinsbaren Theile ist eine Zinsherabsetzung in nicht zu grosser Ferne nicht unwahrscheinlich: jedenfalls aber ist eine Forderung der Vorsicht, bei derartigen Rechnungen immer einen etwas niedrigern Zinsfuss zum Grunde zu legen, als dermalen gangbar ist.

Als Mortalitätstafeln, so weit die Rechnungen davon abhängig sind, habe ich auch diesmal die von \textsc{Brune} benutzt, die zuverlässigsten, die überhaupt vorhanden sind.

Diejenigen Rechnungselemente, welche nur aus den bei der Witwenkasse selbst gemachten Erfahrungen abgeleitet werden können, und also an Zuverlässigkeit gewinnen, wenn diese Erfahrungen einen grössern Zeitraum umfassen, habe ich fur die jetzige Bilanzrechnung sämmtlich neu bestimmt, indem ich die fruhern Erfahrungen mit den neu hinzugekommenen verknüpfte. Bei den einzelnen Positionen wird dies näher angegeben werden.

Endlich bemerke ich noch, dass bei allen Geldangaben Goldwährung zu verstehen ist, und dass die Originalrechnungen zwar durchgehends auf Bruchtheile des Thalers genau gefuhrt, diese Bruchtheile aber in gegenwärtigem Auszuge weggelassen sind. Aus diesem Umstande hat man einige scheinbare kleine Dis%
